\documentclass[10pt,a4paper]{amsart}
\usepackage[margin=19mm]{geometry}
\usepackage{amsmath,amssymb,mathtools}
\usepackage[scr=rsfs,cal=euler]{mathalfa}
\usepackage{xfrac}
\usepackage[unicode,hidelinks,bookmarks=false]{hyperref}
\newcommand{\ie}{i.e.\ }
\newcommand{\cf}{cf.\ }
\newcommand{\resp}{respectively\ }
\newcommand{\Subsection}{\subsection}
\providecommand{\nobreakdash}{\nobreak\hbox{-}}
\setlength{\emergencystretch}{2em}
\multlinegap0pt
\usepackage{amssymb}
\usepackage{mathtools}
\usepackage{braket}
\usepackage{csquotes}
\usepackage{enumerate}
\usepackage{xcolor}
\newtheorem{proposition}{Proposition}
\newcommand{\Ad}{\mathrm{Ad}}
\newcommand{\ddt}{\frac{\de}{\de t}}
\newcommand{\de}{\mathrm{d}}
\newcommand{\g}{\mathfrak{g}}
\title{Sard property for rank 2 polarizations \\ in metabelian Lie groups}
\author{Enrico Le Donne, Antonio Lerario, Luca Nalon, Nicola Paddeu, Luca Rizzi}
\date{Source: arXiv:2607.12530v2 (CC BY 4.0)}
\begin{document}
\maketitle

\noindent\emph{Source and licence.} Sard property for rank 2 polarizations \\ in metabelian Lie groups. Version arXiv:2607.12530v2 (CC BY 4.0), distributed by arXiv under the Creative Commons Attribution 4.0 International licence, http://creativecommons.org/licenses/by/4.0/, as recorded in the arXiv licence metadata for this exact version and shown on the abstract page https://arxiv.org/abs/2607.12530v2. Full licence notice: \texttt{licenses/arxiv-2607.12530-CC-BY-4.0.txt}.\par\smallskip

\noindent\textit{Editorial scope.} Counted unit: the article's proposition proposition (source0.tex lines 269--271). The article's complete original proof of it is delivered with it (source0.tex lines 273--283, one \texttt{proof} environment of 11 lines). No mathematics is written, repaired or re-derived: the statement and the proof are the article's own lines, in the article's own notation, and no retained line is substituted or re-flowed.  No further source-local context is needed: every cross-reference of every retained line resolves inside the two delivered spans.  Nothing was omitted from any retained span, so the package carries no omission record.  No retained line contains a citation, so the wrapper carries no bibliography and no bibliography entry.  No retained line contains a figure environment, an image-inclusion command or drawing code, so no image is delivered and the package declares no asset dependency.  Editorial formatting only, in addition: an AMS \texttt{amsart} preamble that restates the article's own declarations and macros used by the retained lines, namely the environments \texttt{proposition} on separate counters, together with the standard packages those lines need.  The retained environments print in this document as Proposition 1 in order of appearance, whereas the article prints the same items with the numbering of its own full document.  The complete native source of the article as distributed in the arXiv e-print for this exact version is bundled unchanged as \texttt{sources/205123/source0.tex}.
\begin{proposition} \label[proposition]{rem:quotient}
    Let $(G,V)$ be a polarized group with Lie algebra $\g$, and $N \le G$ be a normal Lie subgroup of $G$ with Lie algebra $I \subseteq \g$. Then $(G/N,V/I)$ is a polarized group. Denote by $\pi \colon G \to G/N$ the quotient map and let $\gamma_u$ be an horizontal curve. Then $\pi \circ \gamma_u$ is horizontal for $(G/N,V/I)$ and it is abnormal if and only if $\gamma_u$ is $\lambda$-abnormal for some $\lambda \in \mathbb{S}(I^\perp)$.
\end{proposition}

\begin{proof}
    Clearly, $G/N$ is a Lie group and $V/I= \de\pi(V)$ is bracket-generating whenever $V$ is, since each surjective Lie algebra morphism maps bracket-generating sets in bracket-generating sets. Moreover, if $\gamma_u$ be an horizontal curve in $(G,V)$, then
    \begin{equation*}
        \ddt (\pi \circ \gamma_u) = \de \pi \circ \de L_{\gamma_u(t)}(u(t))= \de L_{\pi(\gamma_u(t))} \de \pi(u(t)) \in  \de L_{\pi(\gamma_u(t))}(\de\pi(V))= \de L_{\pi(\gamma_u(t))}(V/I),
            \end{equation*}
        where in the second equality we used that $\pi$ is a Lie group morphism. For the second part of the statement, we use the canonical identification $\left(\de \pi(\g)\right)^* \cong \ker(\de \pi)^\perp = I^\perp$ and that, since $\pi$ is a Lie group morphism, we have $\Ad_{\pi(g)}(\de \pi(X)) = \de \pi \left(\Ad_g(X)\right)$ for every $g \in G$ and $X \in \g$. Thus
        \begin{equation} \label{poly_proj}
            P^\lambda_X = P^\lambda_{\de\pi(X)} \circ \pi, \quad \text{for every $\lambda \in \mathbb{S}(I^\perp)$.}
        \end{equation}
        Fix $\lambda \in \mathbb{S}(I^\perp)$. On the one hand $\pi \circ \gamma_u$ is $\lambda$-abnormal if and only if $P^\lambda_{\de \pi(X)}(\pi \circ \gamma_u(t))=0$ holds for every $t \in [0,1]$ and $X \in V$. On the other hand, $\gamma_u$ is $\lambda$-abnormal if and only if $P^\lambda_X(\gamma_u(t))=0$ holds for every $t \in [0,1]$ and $X \in V$. In view of \eqref{poly_proj}, the two conditions are equivalent.
\end{proof}

\end{document}
